%
% ---
% id : "electrotechnique_eclairage_20260429_522"
% title : "Calcul du nombre de luminaires LED pour un laboratoire"
% domain : "Électrotechnique"
% subdomain : "Éclairage"
% tags : ["éclairage", "flux lumineux", "éclairement", "rendement", "surface", "led", "ts"]
% difficulty : 1
% ---

\documentclass[a4paper,11pt,answers,addpoints]{exam}
% ---------------------------------------------------------
% LANGUE, ENCODAGE, FONTS
% ---------------------------------------------------------
\usepackage[utf8]{inputenc}

% ---------------------------------------------------------
% GÉOMÉTRIE ET MISE EN PAGE
% ---------------------------------------------------------
\usepackage[left=1.7cm,right=1.5cm,top=2cm,bottom=1cm]{geometry}
%\usepackage{a4wide}
\usepackage{microtype}      % Meilleure répartition du texte
\usepackage{float}          % Positionnement [H]
\usepackage{needspace}      % Saut conditionnel intelligent

% Éviter les veuves/orphelines (optionnel, à décommenter si besoin)
% \widowpenalty=10000
% \clubpenalty=10000

% Espacement vertical flexible
\raggedbottom
\setlength{\parskip}{0.5em plus 0.2em minus 0.1em}
\setlength{\parindent}{0pt}

% ---------------------------------------------------------
% MATHÉMATIQUES
% ---------------------------------------------------------
\usepackage{amsmath, amssymb, bm, esvect}

% Vecteurs
\newcommand{\V}{\overrightarrow}
\def\vect#1{\overrightarrow{\strut#1}}
\providecommand{\Degrees}[0]{\ensuremath{^\circ}}

% ---------------------------------------------------------
% UNITÉS ET GRANDEURS (SIunitx)
% ---------------------------------------------------------
\usepackage{siunitx}
\sisetup{output-decimal-marker={,}}
\DeclareSIUnit \var {var}
\DeclareSIUnit \VA {VA}
\DeclareSIUnit \kVA {kVA}
\DeclareSIUnit[quantity-product={}] \degree{\text{\textdegree}}

% ---------------------------------------------------------
% GRAPHIQUES : TikZ + CircuitikZ + PGFPlots
% ---------------------------------------------------------
\usepackage{tikz}
\usetikzlibrary{
	arrows.meta,
	intersections,
	decorations.pathreplacing,
	positioning
}

\usepackage[european,straightvoltages,EFvoltages]{circuitikz}

\usepackage{tkz-fct}
\usepackage{pgfplots}
\pgfplotsset{compat=1.12}

% Symbole étoile-triangle personnalisé
\newcommand{\wye}{%
	\mathbin{\tikz[x=1ex,y=1ex]{%
			\draw[line width=.1ex] (0,0)--(45:1)--++(-45:1) (45:1)--++(0,1);
	}}%
}

% ---------------------------------------------------------
% TABLEAUX
% ---------------------------------------------------------
\usepackage{tabularx}
\usepackage{tabu}

% ---------------------------------------------------------
% HYPERLIENS
% ---------------------------------------------------------
\usepackage[colorlinks=true,
menucolor=red,
breaklinks=true,
urlcolor=blue,
linkcolor=blue]{hyperref}

% ---------------------------------------------------------
% COMMANDES PERSONNELLES
% ---------------------------------------------------------
\newcommand\nomauteur{bdminasmorgul@protonmail.com}
\newcommand\dateTe{26.04.2026}
\newcommand\brancheTe{TS}

% ---------------------------------------------------------
% STYLE DE L'EXERCICE
% ---------------------------------------------------------
\pagestyle{headandfoot}
\firstpageheadrule
\runningheadrule

\firstpageheader{\brancheTe\ (\nomauteur)}{Élève : }{(\dateTe) \thepage/\numpages}
\runningheader{\brancheTe\ (\nomauteur)}{Élève : }{(\dateTe) \thepage/\numpages}

% Compteur en bas de page
\newcommand{\totalpointtts}{%
	\begin{tikzpicture}[remember picture, overlay]
		\node[anchor=south west, inner sep=0pt, outer sep=0pt]
		at ([xshift=0.5cm,yshift=0.5cm]current page.south west)
		{\rotatebox{90}{Total des points : \numpoints}};
	\end{tikzpicture}
}

\firstpagefooter{\totalpointtts}{}{}
\runningfooter{\totalpointtts}{}{}

% Réglages exam
\printanswers
\addpoints
\pointsinmargin
\bracketedpoints

% ---------------------------------------------------------
% LISTINGS (code)
% ---------------------------------------------------------
\usepackage{xcolor}
\usepackage{listings}

\lstdefinestyle{octavestyle}{
	language=Matlab,
	basicstyle=\ttfamily\small,
	keywordstyle=\color{blue},
	commentstyle=\color{green!50!black},
	stringstyle=\color{red},
	stepnumber=1,
	frame=single,
	breaklines=true,
	breakatwhitespace=true,
	tabsize=4
}

% ---------------------------------------------------------
% DOCUMENT
% ---------------------------------------------------------
\begin{document}
	
	
	\unframedsolutions
	\pointsinmargin
	\bracketedpoints
	\section*{Préparation TS PQ CFC ELMO, IELE, PELE}
	\begin{questions}

\question[9]
Un laboratoire de mesures doit être équipé de nouveaux luminaires LED. Les caractéristiques du projet sont les suivantes : 
\begin{itemize} 
	\item Éclairement moyen requis : $E_m = \qty{750}{[\lux]}$ 
	\item Dimensions du local : Longueur $l = \qty{15}{[\meter]}$, largeur $b = \qty{8}{[\meter]}$ 
	\item Rendement global de l'installation (incluant le facteur de maintenance) : $\eta = \qty{0,6}{}$ 
	\end{itemize} 
	Le luminaire LED choisi possède un flux lumineux $\Phi_L = \qty{8500}{[\lumen]}$. Déterminer le nombre de luminaires n à installer pour respecter ces contraintes~?
\begin{solution} 
	\begin{align*} 
		A &= l \cdot b = \qty{15}{[\meter]} \cdot \qty{8}{[\meter]} = \qty{120}{[\meter^2]} \\[6pt] 
		n &= \dfrac{E_m \cdot A}{\Phi_L \cdot \eta} \\[6pt] 
		n &= \dfrac{\qty{750}{[\lux]} \cdot \qty{120}{[\meter^2]}}{\qty{8500}{[\lumen]} \cdot \qty{0,6}{}} \\[6pt] 
		n &= \dfrac{\qty{90000}{}}{\qty{5100}{}} = \qty{17,65}{} \implies \qty{18}{~luminaires} 
	\end{align*}
	\begin{verbatim} 
		% Télécharger OCTAVE depuis https://wiki.octave.org/Using_Octave et l'installer 
		% Puis copier-coller le code dans OCTAVE
		
		% Données du projet 
		Em = 750;  % Éclairement requis [lux] 
		longueur = 15;      % [m] 
		largeur = 8;        % [m] 
		eta = 0.6;          % Rendement global (facteur maintenance inclus) 
		phi_L = 8500;       % Flux d'un luminaire [lm]
		
		% Calcul de la surface A = longueur * largeur;
		% Calcul du nombre théorique de luminaires 
		n_theo = (Em * A) / (phi_L * eta);
		% Arrondi à l'entier supérieur 
		n_reel = ceil(n_theo);
		% Affichage des résultats 
		printf("Surface à éclairer : %.2f [m^2]\n", A); 
		printf("Nombre théorique de luminaires : %.2f\n", n_theo); 
		printf("Nombre de luminaires à installer : %d\n", n_reel); 
		printf("Éclairement final estimé : %.2f [lux]\n", (n_reel * phi_L * eta) / A); 
	\end{verbatim} 
\end{solution}
\end{questions}
\end{document}
